\chapter{C1 UI: completion of the ordinary/small strip}
\label{component:pa0b854a04f49}
\noindent\textbf{Component }\path{pa0b854a04f49}\par
{\small\noindent Source: \path{R13/baseline/advances/c1_ui/ui_all_tests.tex}\par}
{\small\noindent SHA256: \path{a0b854a04f49c38b64e6b24efb55cafc459729aa6cc777d41b9b5c17041d0e9b}\par}
{\small\noindent Inventory use: CU1--CU13 and entire full-normal/finite-cap proof | C1\_\allowbreak{}BARE\_\allowbreak{}OTHER\_\allowbreak{}TERMINALS; precise subordinate uses retained below\par}
\medskip
\section{C1 UI: completion of the ordinary/small strip}
\label{pa0b854a04f49:sec:r6-c1-ui}
Use the unchanged native C1 source in the representative of (NI7):
$k,\eta$ are positive bounded units and $\zeta=g^2$ is positive
infinitesimal. All original openings and $D=0$ are retained.
The exact stabilizer recovers the other representative; no condition
$k\le1$ is imposed after that change. This section proves the missing
ordinary/small testing strip, and joins the remote theorem (CR1)--(CR13).

\paragraph{Exact bounded remainder.}
Normalize by the positive pure-source $a$. Write $x=k^2$,
$r=1-a\eta$, $K=1+x$, $R=1+xr$, $M=R-t$, and
\[
 J_s=K-xtg^2-a,\quad E=1-x\eta+tx\eta g^2,\quad
 Q_c=1-xg^2r,\quad y=x\eta g^2.
\]
Here $a,t,r,1-r,M,J_s$ are positive bounded openings; $E,Q_c,R,K,r+\eta$
are positive units. The exact polynomials $P_4,P_d,P_\alpha$ are (NC2).
The coefficients after division by $a$ are
\[
\begin{aligned}
 c&=2/a,& A_0&=\frac{2\eta^2P_4}{Kr(r+\eta)E},&
 d_0&=\frac{4\eta P_d}{K(r+\eta)E},\\
 D_0&=\frac{2(1+y)}{J_sE}-d_0/2,&
 \alpha_0&=\frac{2x^2\eta P_\alpha}{KQ_cE},&
 \alpha/a&=g^2\alpha_0.
\end{aligned}                                                     \tag{CU1}
\]
In particular $c,A_0$ are strong; $d_0,\alpha_0$ are units; and
$ca,A_0r,D_0J_s$ are bounded with continuous rational expressions
through every opening. Put
\[
 h_2=\frac{2x\eta(1-tg^2)(1+y)}E,\qquad
 h_b=\frac{2x\eta[1+(1-t)g^2]}E,
 \qquad h_b-h_2=2x\eta g^2.
\]
Both $h_2,h_b$ are positive units. The exact coefficients are
$m/a=h_2/(g^2M)$ and $\beta/a+d_0/2=h_b/(g^2t)$.

For the test put $D_u=1+u^2,D_v=1+v^2$ and
\[
 Q=W+uY,\quad B=X+uZ,\quad p=t/R,
 \quad N=(B-pvQ)/g,\quad L=vQ/g.
\]
The inverse $W=Q-uY$, $X=gN+pvQ-uZ$ is polynomial.
Retain the two positive squares
\[
 \frac{4h_2t}{MR}N^2,\qquad
 \frac{4h_bM}{tR}D_vN^2.                            \tag{CU2}
\]
Their remainder is exactly (CL7). Choose bounded caps
$c_*\le c$, $A_*\le A_0$, $D_*\le D_0$, with $c_*,A_*>0$,
$D_*\ge0$; bounded original weights can be kept unchanged.
Retain the three positive centered excesses
\[
\begin{aligned}
 &4(c-c_*)D_v(r_t-avZ)^2,\quad
 4(A_0-A_*)D_uD_v(Y-rvZ/\eta)^2,\\
 &4(D_0-D_*)(J_t+J_svZ)^2,
\end{aligned}                                                     \tag{CU3}
\]
where
\[
 r_t=uQ-D_uY,\quad J_t=r_t-D_uvZ+ugvN+upv^2Q.
\]
All five retained rows vanish at the actual contact
\[
 (u,v)=(k,k\sqrt\eta g),\quad
 (W,X,Y,Z)=(g/\sqrt\eta,k(tg^2-1),krg/\sqrt\eta,1).
                                                               \tag{CU4}
\]
Set $c_a=2-c_*a$, $A_r=A_0r-A_*r$, $D_J=D_0J_s-D_*J_s$.
After removing (CU2)--(CU3), the entire parent divided by $a$ is
\begin{align*}
G={}&\frac4R\{(h_2+h_bD_v)N^2+2(h_bD_v-h_2)LN
              +[h_2(1-p)+h_bpD_v]L^2\}\\
 &+4c_*D_vr_t^2+8c_aD_vvZr_t-4c_aaD_vv^2Z^2\\
 &+4A_*D_uD_vY^2+\frac{8A_r}{\eta}D_uD_vvZY
                   -\frac{4A_rr}{\eta^2}D_uD_vv^2Z^2\\
 &+4D_*J_t^2-8D_JvZJ_t-4D_JJ_sv^2Z^2\\
 &+4\alpha_0D_u(gvY-gZ)^2
       +2d_0\{J_t^2+D_u(v^2Y^2-2vYZ-v^2Z^2)\}.
                                                               \tag{CU5}
\end{align*}
Thus $G$ is an exact bounded Gram on $(N,Q,Y,vZ,gZ,L)$ whenever
$u$ is a bounded unit and $v$ infinitesimal. Every source denominator
in its coefficients is a unit. Formula (CU5), rather than a scalar
restriction, supplies all normal terms below.

\paragraph{Local contact lifting with all first normals.}
Put $\kappa=k\sqrt\eta$, $\varepsilon=\kappa g$,
$s=R/(k\eta)$, $\delta=u-k$, and define
\[
 E_Q=Q-svZ,\quad E_Y=Y-rvZ/\eta,\quad
 R_h=(vQ-s\varepsilon^2Z)/\varepsilon,
 \quad S=v\delta Z,\quad T=\varepsilon\delta Z.
\]
Let $c_L=4[h_2(1-p)+h_bpD_v]/R$, a positive bounded unit.
The sole potentially unbounded $L$ row can be replaced exactly:
\[
 L=\kappa R_h+kR gZ,\qquad
 L^2-\kappa^2R_h^2=2kR vQZ-k^2R^2g^2Z^2.           \tag{CU6}
\]
Consequently (CU5) equals $c_L\kappa^2R_h^2$
plus $8\kappa(h_bD_v-h_2)NR_h/R$ and a polynomial $G_0$.
Explicitly, in (CU5) replace its first line by
\[
 \frac4R(h_2+h_bD_v)N^2+8k(h_bD_v-h_2)gNZ
                 +c_L(2kRvQZ-k^2R^2g^2Z^2).        \tag{CU7}
\]
All coefficients of $G_0$ are bounded polynomials in $u,v,g$ with
pure-source unit denominators. The coefficient of $NR_h$ is
infinitesimal, since $h_bD_v-h_2=2x\eta g^2+h_bv^2$.

Here are exact prescriptions for every normal coefficient, avoiding
any implicit omission. Substitute $Q=E_Q+svZ$ and
$Y=E_Y+rvZ/\eta$ in $G_0$ and collect the quadratic polynomial
in $(N,E_Q,E_Y)$. The constant coefficient is $\phi(u,\xi)Z^2$,
where $\xi=v^2$ and $\phi\in(g^2,\xi)$. The linear coefficients
of $E_Q,E_Y$ are $vL_Q(u,\xi)Z,vL_Y(u,\xi)Z$; that of $N$ is
$gL_N(u,\xi)Z$, where $L_N\in(g^2,\xi)$. These parity and ideal
assertions follow term by term in (CU5),(CU7). The remaining quadratic
normal block is the full coefficient collection of the same expression.
For explicit pen-and-paper formulas, set
\[
 b=r/\eta,\quad F_0=us-D_ub,\quad
 G_0^{\rm sc}=F_0-D_u+up\xi s,\quad
 H_0=(8D_*+4d_0)G_0^{\rm sc}-8D_J.
\]
Then, writing $D_v=1+\xi$ in these polynomials,
\begin{align*}
 L_Q={}&2c_LkR+8D_vu(c_*F_0+c_a)+u(1+p\xi)H_0,\\
 L_Y={}&-8D_vD_u(c_*F_0+c_a)
       +8D_uD_v(A_*b+A_r/\eta)-D_uH_0\\
      &+8\alpha_0D_ug^2(b\xi-1)+4d_0D_u(b\xi-1),\\
 L_N={}&8k(h_bD_v-h_2)+u\xi H_0.
\end{align*}
The complete scalar coefficient is
\begin{align*}
 \phi={}&c_L(2kRs\xi-k^2R^2g^2)
   +4c_*D_vF_0^2\xi+8c_aD_vF_0\xi-4c_aaD_v\xi\\
 &+4A_*D_uD_vb^2\xi+8(A_r/\eta)D_uD_vb\xi
                    -4A_rrD_uD_v\xi/\eta^2\\
 &+4D_*(G_0^{\rm sc})^2\xi-8D_JG_0^{\rm sc}\xi-4D_JJ_s\xi\\
 &+4\alpha_0D_ug^2(b\xi-1)^2
   +2d_0\{(G_0^{\rm sc})^2\xi+D_u(b^2\xi^2-2b\xi-\xi)\}.
\end{align*}
These formulas display the two asserted ideals directly. They also
give a short explicit input for the monic divisions, without an
unprinted state-normal computation.

The source is an exact stationary zero and every removed square
vanishes there. Therefore
\[
 \phi=\phi_u=\phi_\xi=0,\qquad L_Q=L_Y=L_N=0
                       \quad\text{at }(u,\xi)=(k,\varepsilon^2).
\]
Successive monic polynomial division gives, with bounded coefficients,
\[
 \phi=\delta^2(\varepsilon^2B_0+\xi B_1)
       +\delta(\xi-\varepsilon^2)B_2
       +(\xi-\varepsilon^2)^2B_3,                  \tag{CU8}
\]
\[
 L_i=\delta L_{i0}+(\xi-\varepsilon^2)L_{i1} (i=Q,Y),\quad
 L_N=\delta(g^2L_{N0}+\xi L_{N1})
                         +(\xi-\varepsilon^2)L_{N2}.
\]
For clarity, each quotient is obtained by subtracting the value at
$u=k$ and dividing by the monic $u-k$, then by the monic
$\xi-\varepsilon^2$. The second $u$ quotient of $\phi$ stays in
$(g^2,\xi)$; its constant and first $u$ terms have the stated double
and simple zeros in $\xi$. The same argument gives the $L_N$ ideal.
No small source number is inverted in these divisions.

Put $R_s=(v^2-\varepsilon^2)Z/\varepsilon$. The exact root identities
\[
 \varepsilon R_s=(\varepsilon R_h-vE_Q)/s,\quad
 R_sT=(R_hT-SE_Q)/s,\quad
 S^2-T^2=\delta(R_hT-SE_Q)/s                       \tag{CU9}
\]
allocate (CU8) and every first normal on the six roots
$(N,E_Q,E_Y,R_h,S,T)$. For example $v\delta Z=S$,
$\delta(\xi-\varepsilon^2)Z^2=(R_hT-SE_Q)/s$, and
$(\xi-\varepsilon^2)^2Z^2=(\varepsilon R_h-vE_Q)^2/s^2$.
The $N$ normal has only infinitesimal couplings. All quadratic normal
coefficients are already bounded. This is the required exact full lift.

\paragraph{Strict matched Hessian and source boundary caps.}
First take finite interior source residues. The leading form is
(CL2). Put $V=k\sqrt\eta V_1$, $W=W_1/\sqrt\eta$,
$Y=kY_1/\sqrt\eta$, $Z=1$. At the contact
$(W_1,Y_1,V_1,u)=(1,r,1,k)$, the Schur complement of the two state
coordinates in half the Hessian, in coordinates $(V_1,u)$, is
\[
\begin{aligned}
 G_{11}&=\frac{32\eta x^2P_4}{KE_0},&
 G_{12}&=\frac{16\eta k^3(R^2-2KE_0)}{KE_0R},\\
 G_{22}&=\frac{32\eta x(P_4+xrP_d)}{R^2P_4},&
 \det G&=\frac{256\eta^2x^3(4K^2E_0-R^2)}{K^2E_0^2}.
\end{aligned}                                                     \tag{CU10}
\]
Here $E_0=1-x\eta$, $P_4=KE_0+x(1-r)$, $P_d=1+2x+rx^2$.
There are no opening denominators. The leading $D\ge0$ condition is
\[
 K(2+xr)E_0-R^2\ge0,
\]
so $4K^2E_0-R^2\ge K(4K-2-xr)E_0>0$.
The state block is uniformly bounded below since $c,A_0$ are strong.
Thus the matched Hessian is strict, uniformly through every permitted
source boundary after the appropriate hard directions are retained.

We give the boundary justification explicitly. Divergences can occur
only at $a=0$, $r=0$, or $J_s=0$, for $c,A_0,D_0$, respectively.
Since $a+J_s\sim K$, $c$ and $D_0$ cannot both diverge. If $a=0$,
then $r=1$, so $c$ is the only divergent weight. The only possible
divergent pair is $A_0,D_0$. Their hard-state rows are $E_Y$ and
$kE_Q-KE_Y$, with nonzero determinant $k$. All angular parts of
these centered rows are bounded. Solving these at most two real
linear equations therefore supplies a bounded competitor for the
minimizing state vector. The bounded finite part and (CU10) imply
that the state minimizer is bounded as well, by completing its
positive state square. The full matched Hessian hence has a uniform
positive lower bound on the Euclidean norm in a nearby finite-source
family. On a vector in the limiting hard kernel, adjust only the two
state coordinates to annihilate the nearby exact hard rows. This
adjustment tends to zero, and the divergent penalties vanish exactly.
The bounded finite part is consequently strictly positive on that
kernel. The elementary real finite-penalty lemma supplies finite caps.

Every allowed residue is approachable by such finite sources. For
fixed $x>0$, the real leading inequalities are
\[
 0\le r\le1,\qquad
 \eta_{\min}(r)=\frac{1-r}{K}\le\eta\le
 \eta_{\max}(r)=\frac{1-R^2/[K(2+xr)]}{x},
\]
with positive interval width $R/[xK(2+xr)]$. Perturb $r$ into
$(0,1)$ and its interval coordinate into $(0,1)$.
This makes $a,r,J$ and the $D$ numerator strictly positive while
converging to the given residue. Taking $t=R/2$ and then sufficiently
small real $g>0$ restores all original source signs, including
$\beta>0$. This construction, rather than substitution into a
singular leading denominator, justifies the boundary limit.

In the six-root residue Gram, $N$ is an isolated positive pivot.
The $(E_Q,E_Y)$ block is strict, $R_h$ has positive pivot
$c_L\kappa^2$, $S$ couples only to $E_Q,E_Y$, and $T$ only to $R_h$.
Setting $S=T$ is the matched Hessian under an invertible real change
of coordinates: $R_h/g$ has relative-angle derivative $2s\kappa\ne0$,
and $T/g=\kappa\delta$ supplies the ordinary-angle derivative.
After the hard Schur complements, the sum of the $S,T$ diagonals is
therefore strictly positive. Add a real multiple of the last zero
in (CU9) to redistribute this sum into two positive diagonals.
The extra coefficients contain $\delta$ and vanish in the residue.
The bounded strict-residue Gram lemma proves the full parent at
$u-k$ and $v$ infinitesimal, including every source opening.

\paragraph{Ordinary angles away from contact.}
Let now $\bar u^2\ne\bar k^2$, with $u$ a unit and $v$ infinitesimal.
Use the bounded frame $(N,r_t,Y,vZ,gZ,L)$ of (CU5), and the exact zero
\[
 L(gZ)-Q(vZ)=0,\qquad Q=(r_t+D_uY)/u.              \tag{CU11}
\]
For finite leading source weights, $N$ is an isolated positive pivot.
The remaining residue is
\[
 q^TPq+2Ub^Tq+c_2U^2+4hL^2+c_0T^2,
 \quad q=(r_t,Y)^T,\ U=vZ,\ T=gZ,
\]
where the source coefficients are (CL1) and
\[
\begin{gathered}
 s_D=4D+2d,\quad s_C=4C+s_D,\quad
 P=\operatorname{diag}(s_C,4AD_u),\\
 \ell=(1/u,D_u/u)^T,\quad b=(-s_DD_u,-2dD_u)^T,\\
 c_0=4\alpha D_u,\qquad c_2=s_DD_u^2-2dD_u.
\end{gathered}                                                     \tag{CU12}
\]
Apply the real rank-one allocation lemma (UA1)--(UA3), using (CU11).
We verify its hypotheses uniformly through all source boundaries.
Writing $d_a=\ell^TP^{-1}\ell$, $e_a=\ell^TP^{-1}b$,
$h_a=c_2-b^TP^{-1}b$, direct inversion yields
\[
\begin{aligned}
 d_a&=\frac{aJE_0}{8u^2(JE_0+a)}+\frac{D_u}{4Au^2},\\
 e_a&=-\frac{aD_u}{u(JE_0+a)}-\frac{dD_u}{2Au},\\
 h_a&=\frac{8D_u^2}{JE_0+a}-2dD_u-\frac{d^2D_u}{A},\\
 P^{-1}b&=(-aD_u/(JE_0+a),-d/(2A))^T.
\end{aligned}                                                     \tag{CU13}
\]
All are bounded continuous, since $JE_0+a$ is a positive unit and
$1/A$ has a continuous bounded formula proportional to $r$.
The same holds for $P^{-1}\ell$.

Normalize the Laurent polynomial (CL4) by $As_C$.
Its $U$ and $T^{-1}$ coefficients are, respectively,
$1/[2(JE_0+a)]$ and $\alpha/4$, positive units.
Its $T$ coefficient is
\[
 h\{(D+d/2)/s_C+2CD/(As_C)+d/(8A)\},
\]
also a positive unit: when $1/A=0$, $r=0$ forces $a$ a unit,
so $(D+d/2)/s_C=a/[4(JE_0+a)]$ is a unit; otherwise the last term
is positive. The $U^{-1}$ coefficient is
$h\alpha(1/s_C+1/(4A))$. If it vanishes, $r=0$ and $J=0$,
and the $T/U$ coefficient has the positive term $hD/s_C=h/4$.
All remaining normalized nonconstant coefficients are bounded and
nonnegative. The constant is bounded by the exact contact equation.
Thus strict convexity and coercivity in logarithmic coordinates
persist at every boundary. The normalized Laurent function retains
its unique stationary zero $(U,T)=(x,x\eta)$.

For fixed $u^2\ne x$, the rank-one determinant polynomial is
\[
 p(T)=c_0+(h_a+4hd_ac_0)T+4h(d_ah_a+e_a^2)T^2
      =\frac{16D_u}{As_C}T\mathcal L(u^2,T)>0.
\]
Its constant and leading coefficients are positive units by the
preceding coefficient bounds. The strict allocation interval
therefore survives at every source boundary. The two real Gram
blocks have a uniform positive lower bound; (CU13) controls their
state minimizers even if $P$ is unbounded. The centered divergent
rows have hard parts $r_t$, $Y$, $r_t$; the only simultaneous pair
is $Y,r_t$, whose determinant is one. Exact small state corrections
annihilate the large rows in the nearby finite family. The same
finite-penalty argument as above supplies caps in (CU3). The exact
bounded remainder (CU5), after (CU11), has strict residue.
This proves every ordinary/small test away from contact.

\paragraph{Exhaustive admission.}
The two arguments cover the whole ordinary/small testing strip;
the exact gauges cover negative contact signs. The remote theorem
(CR1)--(CR13), including CR5a, covers its complement.
All transformations are homogeneous in the state; no state coordinate
was divided out in the tensor proof. Real and imaginary state parts
are treated separately. Hence the entire actual UI subdomain of C1.5,
and its exact stabilizer images, has an all-tests tensor-SOS proof.
The actual UH subdomain of C1.5 remains a separate obligation.
